\documentclass[10pt,a4paper]{amsart}
\usepackage[margin=19mm]{geometry}
\usepackage{amsmath,amssymb,mathtools}
\usepackage[scr=rsfs,cal=euler]{mathalfa}
\usepackage{xfrac}
\usepackage[unicode,hidelinks,bookmarks=false]{hyperref}
\newcommand{\ie}{i.e.\ }
\newcommand{\cf}{cf.\ }
\newcommand{\resp}{respectively\ }
\newcommand{\Subsection}{\subsection}
\providecommand{\nobreakdash}{\nobreak\hbox{-}}
\setlength{\emergencystretch}{2em}
\multlinegap0pt
\usepackage{amssymb}
\usepackage{xcolor}
\usepackage{dsfont}
\newtheorem{thm}{Theorem}
\newtheorem{prop}{Proposition}
\newcommand{\DP}{\mathds{D}\mathds{P}}
\newcommand{\I}{\mathcal{I}}
\newcommand{\J}{\mathcal{J}}
\newcommand{\N}{\mathds{N}}
\newcommand{\R}{\mathds{R}}
\newcommand{\Z}{\mathds{Z}}
\newcommand{\de}{\partial}
\newcommand{\pk}{para-K\"ahler }
\newcommand{\s}{\mathcal{S}}
\title{Para-K\"ahler immersions in \pk space forms}
\author{Gianni Manno, Filippo Salis}
\date{Source: arXiv:2503.11582v1 (CC BY 4.0)}
\begin{document}
\maketitle

\noindent\emph{Source and licence.} Para-K\"ahler immersions in \pk space forms. Version arXiv:2503.11582v1 (CC BY 4.0), distributed by arXiv under the Creative Commons Attribution 4.0 International licence, http://creativecommons.org/licenses/by/4.0/, as recorded in the arXiv licence metadata for this exact version and shown on the abstract page https://arxiv.org/abs/2503.11582v1. Full licence notice: \texttt{licenses/arxiv-2503.11582-CC-BY-4.0.txt}.\par\smallskip

\noindent\textit{Editorial scope.} Counted unit: the article's thm thm (source0.tex lines 905--908). The article's complete original proof of it is delivered with it (source0.tex lines 909--983, one \texttt{proof} environment of 75 lines). No mathematics is written, repaired or re-derived: the statement and the proof are the article's own lines, in the article's own notation, and no retained line is substituted or re-flowed.  Retained uncounted as the source-local context the proof needs: lines 206--206; lines 279--280; lines 295--296; lines 339--343; lines 423--429; lines 431--433; lines 438--442; lines 456--479; lines 465--468.  Nothing was omitted from any retained span, so the package carries no omission record.  No retained line contains a citation, so the wrapper carries no bibliography and no bibliography entry.  No retained line contains a figure environment, an image-inclusion command or drawing code, so no image is delivered and the package declares no asset dependency.  Editorial formatting only, in addition: an AMS \texttt{amsart} preamble that restates the article's own declarations and macros used by the retained lines, namely the environments \texttt{thm}, \texttt{prop} on separate counters, together with the standard packages those lines need.  The retained environments print in this document as Theorem 1, Proposition 1 in order of appearance, whereas the article prints the same items with the numbering of its own full document.  The complete native source of the article as distributed in the arXiv e-print for this exact version is bundled unchanged as \texttt{sources/174772/source0.tex}.
\subsection{Para-K\"ahler  manifolds, diastasis and \pk space forms}\label{sec.pk.diastasis}

\begin{equation}\label{dn}
 \Phi=4\sum_{i=1}^N \xi_i\eta_i =D_0(\xi,\eta), \end{equation}

\begin{equation}\label{dpn}
\Phi= \frac{8}{c}\log\left( 1+2\sum_{i=1}^N \xi_i\eta_i\right)=D_0(\xi,\eta). \end{equation}

\begin{thm}[Rigidity]\label{rigidity}
Let $f$ and $g$ be  two weakly full \pk immersions from an open subset $U$ of a \pk manifold $(M^n,\omega)$ into $(\s_c^{N_1},\omega_c)$ and  $(\s_c^{N_2},\omega_c)$, respectively. Let also assume that
the diastasis function is defined on $U\times U$.  Then,  $N_1=N_2$ and 
$f$ differs from $g$ for a para-holomorphic isometry of the ambient space.
\end{thm}

\begin{equation}\label{H}
H_c(\xi,\eta):=
\begin{cases}
\frac{1}{4}D_0^{M^n}(\xi e+\eta \bar e)&\text{if } c=0;\\
\frac{1}{2}\exp\left( \frac{c}{8}D_0^{M^n}(\xi e+\eta \bar e)\right)-\frac{1}{2} &\text{if } c\neq 0,\\
\end{cases}
\end{equation}

\begin{equation}\label{problem}
H_c(\xi,\eta)=\sum_{\alpha=1}^N u_\alpha(\xi)v_\alpha(\eta).
\end{equation}

\begin{prop}\label{existence}
Let $U$ be an open subset of \pk manifold $(M^n,\omega)$ where it is defined a system of null-coordinates $(\xi,\eta)$. Then, there exists a weakly full \pk immersion
 $f: (U,\omega|_U)\to (\s_c^N,\omega_c)$
if and only if the function $H_c$ defined by \eqref{H} reads as \eqref{problem} with the smallest possible $N$. 
\end{prop}

\begin{thm}\label{main}
Let $p$ be a point of a \pk manifold $(M^n,\omega)$. Let $U=\Omega\times\Omega$ be a neighborhood of $p$ where it is defined a system of null-coordinates $(\xi,\eta)$ centered at $p$. Let the diastasis function be defined on $U\times U$.
Then  necessary conditions for the existence of a  weakly  full \pk immersion 
$$f: (U,\omega|_U)\to (\s_c^N,\omega_c)$$
are the following:\\
There exist two sets  $\I,\J\subset \N^n$ of multi-indices, containing 0 and having finite cardinality, for which $\forall\, K\in\N^n$ there exist
smooth functions $a^I_K, b^J_K:\Omega \to\R $, where $I\in\I$, $J\in\J$, $K\in\N^n$ and two  smooth functions 
$a, b$ not identically zero in a neighborhood  of $0$ with $a(0)=b(0)=0$, 
such that
\begin{equation}\label{dip1}
a(\xi)\ \frac{\de^{|K|}H_c}{{\de\xi^{K}}}(\xi,\eta) + \sum_{I\in\I} a^I_K(\xi)\   \frac{\de^{|I|}H_c}{\de\xi^{I}} (\xi,\eta) \equiv 0 
%,\qquad\forall K\in\N^n,
\end{equation}
and
\begin{equation}\label{dip2}
b(\eta)\ \frac{\de^{|K|}H_c}{{\de\eta^{K}}}(\xi,\eta) + \sum_{J\in\J} b^J_K(\eta)\   \frac{\de^{|J|}H_c}{\de\eta^{J}} (\xi,\eta) \equiv 0
%,\qquad \forall K\in\N^n,
\end{equation}
where the function $H_c$ is defined by \eqref{H}.  The above conditions are sufficient by assuming that either \begin{equation}\label{eq:A}
\Omega'=\big\{\xi\in\Omega\ | \ a(\xi)\neq 0\big\}\,,
\end{equation}
or $\big\{\eta\in\Omega\ | \ b(\eta)\neq 0\big\}$,
is a connected  open dense subset of $\Omega$. This leads actually to a strongly full \pk immersion.
\end{thm}

\begin{equation}
a(\xi)\ \frac{\de^{|K|}H_c}{{\de\xi^{K}}}(\xi,\eta) + \sum_{I\in\I} a^I_K(\xi)\   \frac{\de^{|I|}H_c}{\de\xi^{I}} (\xi,\eta) \equiv 0 
%,\qquad\forall K\in\N^n,
\end{equation}

\begin{thm}
A \pk space form $(\s_c^n,\omega_c)$ can be locally \pk immersed into $(\s_b^N,\omega_b)$, where $N\geq n$, if and only if either $c=b=0$ or $b/c\in\Z^+$. Such local immersions can be extended to the whole manifold. Moreover, if
$N=n$ (when $c=b=0$), or if $N =\binom{n+\frac{b}{c}}{n}-1$ (when $\frac{b}{c}\in\Z^+$), these immersions are strongly full.
\end{thm}

\begin{proof}

Since \pk space forms are homogeneous with respect to the action of their para-homolorphic isometry group (see the end of Section \ref{sec.pk.diastasis}), the existence (or non-existence) of  a \pk immersion defined on a neighborhood of an arbitrary chosen point implies the existence (or non-existence) of a \pk immersion defined on a neighborhood of any other point. Therefore, we can study, without loss of generality, only local \pk immersions defined either on a neighborhood of  $0\in \s_0^n$ or on a neighborhood of  $[1,0\mathellipsis,0]\in\s^n_c$, $c\neq 0$. Let us fix a null coordinate system $(\xi,\eta)$ around such point. Keeping in mind the definitions \eqref{dn}, \eqref{dpn} and \eqref{H},
we compute the function $H_b$ in order to apply Theorem \ref{main}.
Notice that a flat complex space form $(\s_0^n,\omega_0)$ can be trivially \pk immersed by inclusion into $(\s_0^N,\omega_0)$, hence only the following cases will be taken into account:
\begin{itemize} 
\item $c\neq 0$ and $b= 0$;
\item $c= 0$ and $b\neq 0$;
\item $c\neq 0$ and $b\neq 0$.
\end{itemize} 
Below we treat the above cases separately.


\medskip
\noindent
$\bullet$ If $c\neq 0$ and $b= 0$, by taking into account   \eqref{H} and that the diastasis $D_0^{\s_c^n}$ is given by \eqref{dpn}, we have that
$$H_0(\xi,\eta)= \frac{2}{c}\log\left( 1+2\sum_{i=1}^n \xi_i\eta_i\right).$$
By differentiating the previous equation with respect to  $\xi_1$, we get, for any $k\in\Z^+$,
$$\frac{\de^k H_0}{\de \xi_1^{k}} = \frac{(-2)^{k+1}(k-1)!}{c \left( 1+2\sum_{i=1}^n \xi_i\eta_i\right)^k}\ \eta_1^k\,.$$
Let now assume that the condition \eqref{dip1} of Theorem \ref{main} holds true for any  $\I\subset \N^n$ such that $\#\I=N+1$. In particular, if 
\begin{equation}\label{II}
\I=\{(k,0,\mathellipsis,0)\ |\ k=1,\mathellipsis,N+1\},\end{equation}
we have that
$$\sum_{k=1}^{N+1} a_k(\xi) \frac{\de^k H_0}{\de \xi_1^{k}} (\xi,\eta)=\frac{ \sum_{k=1}^{N+1} (-2)^{k+1}(k-1)!  \left( 1+2\sum_{i=1}^n \xi_i\eta_i\right)^{N+1-k} a_k(\xi) \eta_1^k}{c \left( 1+2\sum_{i=1}^n \xi_i\eta_i\right)^{N+1}} \equiv 0.$$
%Since
%$$\sum_{k=1}^N (-2)^{k+1}(k-1)!  \left( 1+2\sum_{i=1}^n \xi_i\eta_i\right)^{N-k} a_k(\xi) \eta_1^k=
%4\left( 1+2\sum_{i=1}^n \xi_i\eta_i\right)^{N-1}a_1(\xi)  \eta_1+\sum_{k=2}^{N} p_k(\xi,\eta) \eta_1^k$$
%
%\textbf{GIANNI: SECONDO ME:}
%$$\sum_{k=1}^N (-2)^{k+1}(k-1)!  \left( 1+2\sum_{i=1}^n \xi_i\eta_i\right)^{N-k} a_k(\xi) \eta_1^k=
%4\left( 1+2\sum_{i=2}^n \xi_i\eta_i\right)^{N-1}a_1(\xi)  \eta_1+\sum_{k=2}^{N} p_k(\xi,\eta) \eta_1^k$$
%
%for some suitable smooth functions $p_k$ independent of $\eta_1$, then $a_1(\xi)\equiv 0$. 
%
%\textbf{E' UN POLINOMIO IN $\eta_1$ IL CUI TERMINE DI GRADO 1 E' $4a_1(\xi)$}
Since the numerator in the above equality is a polynomial in $\eta_1,\mathellipsis,\eta_n$ and the coefficient of its monomial in $\eta_1$ of degree 1 is $4a_1(\xi)$, it follows that $a_1(\xi)$ needs to be identically zero.
 By means of similar considerations, we get also that $a_k(\xi)\equiv 0$ for any  $k=2,\mathellipsis,N+1$.
In view of  Theorem  \ref{main}, we  conclude that there  are no local  \pk immersions of a non-flat \pk space form into a flat one.\\


\medskip
\noindent
$\bullet$ 
If $c=0$ and $b\neq 0$, by taking into account   \eqref{H} and that the diastasis $D_0^{\s_0^n}$ is given by \eqref{dn}, we have that
$$ H_b(\xi,\eta)= \frac{1}{2}\exp\left( \frac{b}{2}\sum_{i=1}^n \xi_i\eta_i \right)-\frac{1}{2}.$$
By differentiating the previous equation with respect to  $\xi_1$, we get, for any $k\in\Z^+$,
$$\frac{\de^k H_b}{\de \xi_1^{k}}=\frac{b^k}{2^{k+1}}\exp\left( \frac{b}{2}\sum_{i=1}^n \xi_i\eta_i \right)\ \eta_1^k\,. $$
Now we assume that the condition \eqref{dip1} of Theorem \ref{main} holds true for any  $\I\subset \N^n$ such that $\#\I=N+1$. In particular, if we consider the subset $\I$ of multi-indices \eqref{II}, we have that
$$\sum_{k=1}^{N+1} a_k(\xi) \frac{\de^k H_b}{\de \xi_1^{k}}(\xi,\eta)= \exp\left( \frac{b}{2}\sum_{i=1}^n \xi_i\eta_i \right) \sum_{k=1}^{N+1}\frac{b^k  a_k(\xi) }{2^{k+1}} \eta_1^k \equiv 0,$$
that implies $a_k(\xi)\equiv 0$ for any  $k=1,\mathellipsis,N+1$.
In view of  Theorem  \ref{main}, we  conclude that there are no local \pk immersions of a flat \pk space forms into a non-flat one.\\


\medskip
\noindent
$\bullet$ 
If $c\neq 0$ and $b\neq 0$, by taking into account   \eqref{H} and that the diastasis $D_0^{\s_c^n}$ is given by \eqref{dpn}, we have that
$$ H_b(\xi,\eta)= \frac{1}{2}\left(1+2\sum_{i=1}^n \xi_i\eta_i\right)^\frac{b}{c}-\frac{1}{2}.$$
By differentiating the previous equation with respect to  $\xi_1$, we get, for any $k\in\Z^+$,
$$\frac{\de^k}{\de \xi_1^{k}} H_b= 
2^{k-1} \prod_{j=0}^{k-1}\left(\frac{b}{c}-j\right)\left(1+2\sum_{i=1}^n \xi_i\eta_i\right)^{\frac{b}{c} -k}\eta_1^k\,.$$
Now we assume that the condition \eqref{dip1} of Theorem \ref{main} holds true for any  $\I\subset \N^n$ such that $\#\I=N+1$. In particular, if we consider the subset $\I$ of multi-indices \eqref{II}, we have that
$$\sum_{k=1}^{N+1} a_k(\xi) \frac{\de^k H_b}{\de \xi_1^{k}}(\xi,\eta)=\left(1+2\sum_{i=1}^n \xi_i\eta_i\right)^{\frac{b}{c} -N-1}
\sum_{k=1}^{N+1} a_k(\xi) 2^{k-1} \prod_{j=0}^{k-1}\left(\frac{b}{c}-j\right)\left(1+2\sum_{i=1}^n \xi_i\eta_i\right)^{N+1 -k}\eta_1^k \equiv 0.$$
%Since
%$$\sum_{k=1}^N a_k(\xi) 2^{k-1} \prod_{j=0}^{k-1}\left(\frac{b}{c}-j\right)\left(1+2\sum_{i=1}^n \xi_i\eta_i\right)^{N -k}\eta_1^k = \left(1+2\sum_{i=2}^n \xi_i\eta_i\right)^{N -1} a_1(\xi) \eta_1+ \sum_{k=2}^N q_k(\xi,\eta)\eta_1^k $$
%for some smooth functions $q_k$, we have $a_1(\xi)\equiv 0$.
Since the above equality can be seen as a polynomial in $\eta_1,\mathellipsis,\eta_n$ and the coefficient of its monomial in $\eta_1$ of degree 1 is $\frac{b}{c}a_1(\xi)$, it follows that $a_1(\xi)$ needs to be identically zero.
If $\frac{b}{c}\not\in\Z^+$, then we get, by means of similar considerations, that also $a_k(\xi)\equiv 0$ for any  $k=2,\mathellipsis,N+1$. Otherwise, if $\frac{b}{c}\in\Z^+$, then $H_b$ reads as \eqref{problem}, so, in view of Proposition \ref{existence} and Theorem \ref{main}, there exists, for any $i=0,\mathellipsis,n$,
a strongly full local \pk immersion 
$$f_i:\big( \mathcal{U}_i=\{[Z]\in\DP^n\ | \ |Z_i|^2\neq0\},\omega_c\big) \to \left(\s_b^{\binom{n+\frac{b}{c}}{n}-1},\omega_b\right)\,.$$
Indeed, we need at least $\binom{n+\frac{b}{c}}{n}-1$ monomials to describe the polynomial $H_b$.

Finally, we notice that it follows from the Rigidity Theorem \ref{rigidity} the existence of a para-holomorphic isometry  $F$ of the ambient space such that  $f_i|_{\mathcal{U}_i\cap\mathcal{U}_j }= F\circ f_j|_{\mathcal{U}_i\cap\mathcal{U}_j }$. Hence, we have that any local immersion $f_i$ can be extended to a global immersion.
\end{proof}

\end{document}
